\section{Requerimientos}

Dentro del proceso de desarrollo interno de software en TenarisTamsa, las especificaciones que el sistema debe cumplir (es decir, los requerimientos de software) son definidas por personal de la propia compañía. Dichas especificaciones establecen las necesidades y funcionalidades que deben contemplarse para garantizar el cumplimiento integral de los objetivos del sistema.

A través del equipo de desarrollo se lleva a cabo un proceso de análisis y negociación de estas propuestas, con el fin de alcanzar un consenso claro y específico sobre las capacidades, alcances y limitaciones que tendrá el sistema.

El Sistema Integral de Producción está conformado por distintos módulos internos, cada uno orientado a cubrir las necesidades de diversas áreas operativas de la corporación. Debido a ello, diferentes grupos de trabajo pueden establecer consensos particulares respecto a las funcionalidades y características que debe incluir cada módulo antes de su implementación dentro de las áreas productivas de la empresa.

No obstante, es importante destacar que todos los módulos pertenecientes al Sistema Integral de Producción deben cumplir con una serie de requerimientos generales establecidos conforme a los estándares corporativos definidos por Tenaris.

\subsection{Requerimientos funcionales}

\textbf{Inicio de sesión}: El sistema deberá permitir el inicio de sesión únicamente a los usuarios que hayan sido previamente registrados en el módulo correspondiente por un usuario con permisos de administrador o gerente (Tabla \ref{tab:req_sesion}).

\begin{table}[H]
    \centering
    \caption{Requerimiento funcional - Inicio de sesión}
    \label{tab:req_sesion}
    \begin{tabularx}{\textwidth}{|X|X|}
        \hline
        Código de requerimiento & RF-01                                                                                                                                                                                                                                             \\ \hline
        Nombre                  & Inicio de sesión                                                                                                                                                                                                                                  \\ \hline
        Propósito               & Permitir el acceso seguro al sistema únicamente a usuarios autorizados.                                                                                                                                                                           \\ \hline
        Descripción general     & El sistema deberá permitir el inicio de sesión a los usuarios previamente registrados en el módulo correspondiente mediante el uso de credenciales válidas, asignando los permisos de acceso de acuerdo con el rol establecido para cada usuario. \\ \hline
        Prioridad               & Alta                                                                                                                                                                                                                                              \\ \hline
    \end{tabularx}
\end{table}


\textbf{Gestión de usuarios}: El sistema deberá permitir a los usuarios con rol de administrador o gerente gestionar las cuentas de usuario del módulo, incluyendo la creación, modificación y eliminación de usuarios (Tabla \ref{tab:req_usuarios}).

\begin{table}[H]
    \centering
    \caption{Requerimiento funcional - Gestión de usuarios}
    \label{tab:req_usuarios}
    \begin{tabularx}{\textwidth}{|X|X|}
        \hline
        Código de requerimiento & RF-02                                                                                                                                                                                                                         \\ \hline
        Nombre                  & Gestión de usuarios                                                                                                                                                                                                           \\ \hline
        Propósito               & Administrar las cuentas de usuario y sus permisos dentro del sistema.                                                                                                                                                         \\ \hline
        Descripción general     & El sistema deberá permitir a los usuarios con rol de administrador o gerente crear, modificar y eliminar cuentas de usuario dentro del módulo correspondiente, así como asignar los permisos de acceso según el rol definido. \\ \hline
        Prioridad               & Alta                                                                                                                                                                                                                          \\ \hline
    \end{tabularx}
\end{table}


\textbf{Reporte de producción}: El sistema deberá permitir a los usuarios generar reportes de producción dentro del módulo correspondiente (Tabla \ref{tab:req_reporte}).

\begin{table}[H]
    \centering
    \caption{Requerimiento funcional - Reporte de producción}
    \label{tab:req_reporte}
    \begin{tabularx}{\textwidth}{|X|X|}
        \hline
        Código de requerimiento & RF-03                                                                                                                                                                                       \\ \hline
        Nombre                  & Reporte de producción                                                                                                                                                                       \\ \hline
        Propósito               & Registrar la información relacionada con la producción realizada.                                                                                                                           \\ \hline
        Descripción general     & El sistema deberá permitir a los usuarios generar reportes de producción mediante el registro de la información correspondiente a las actividades realizadas durante el proceso productivo. \\ \hline
        Prioridad               & Alta                                                                                                                                                                                        \\ \hline
    \end{tabularx}
\end{table}


\textbf{Gestión de registros de producción}: El sistema deberá permitir a los usuarios con rol de operador modificar o eliminar únicamente sus propios registros de producción. Asimismo, los usuarios con rol de supervisor o administrador podrán modificar o eliminar los registros de producción de cualquier usuario (Tabla \ref{tab:req_registros}).

\begin{table}[H]
    \centering
    \caption{Requerimiento funcional - Gestión de registros de producción}
    \label{tab:req_registros}
    \begin{tabularx}{\textwidth}{|X|X|}
        \hline
        Código de requerimiento & RF-04                                                                                                                                                                                                                                                                         \\ \hline
        Nombre                  & Gestión de registros de producción                                                                                                                                                                                                                                            \\ \hline
        Propósito               & Mantener actualizada y controlada la información de producción registrada en el sistema.                                                                                                                                                                                      \\ \hline
        Descripción general     & El sistema deberá permitir a los usuarios con rol de operador modificar o eliminar únicamente sus propios registros de producción. Asimismo, los usuarios con rol de supervisor o administrador podrán modificar o eliminar los registros de producción de cualquier usuario. \\ \hline
        Prioridad               & Alta                                                                                                                                                                                                                                                                          \\ \hline
    \end{tabularx}
\end{table}


\textbf{Registro de almacén}: El sistema deberá permitir a los usuarios con rol de supervisor crear, modificar y eliminar registros relacionados con las materias primas disponibles en almacén (Tabla \ref{tab:req_almacen}).

\begin{table}[H]
    \centering
    \caption{Requerimiento funcional - Registro de almacén}
    \label{tab:req_almacen}
    \begin{tabularx}{\textwidth}{|X|X|}
        \hline
        Código de requerimiento & RF-05                                                                                                                                                             \\ \hline
        Nombre                  & Registro de almacén                                                                                                                                               \\ \hline
        Propósito               & Gestionar la información relacionada con las materias primas disponibles en almacén.                                                                              \\ \hline
        Descripción general     & El sistema deberá permitir a los usuarios con rol de supervisor crear, modificar y eliminar registros relacionados con las materias primas existentes en almacén. \\ \hline
        Prioridad               & Media                                                                                                                                                             \\ \hline
    \end{tabularx}
\end{table}


\textbf{Visualización de la producción}: El sistema deberá permitir a los usuarios con rol de supervisor consultar la información de producción correspondiente al día y turno seleccionados (Tabla \ref{tab:req_visualizacion}).

\begin{table}[H]
    \centering
    \caption{Requerimiento funcional - Visualización de la producción}
    \label{tab:req_visualizacion}
    \begin{tabularx}{\textwidth}{|X|X|}
        \hline
        Código de requerimiento & RF-06                                                                                                                                                 \\ \hline
        Nombre                  & Visualización de la producción                                                                                                                        \\ \hline
        Propósito               & Consultar la información de producción registrada en el sistema.                                                                                      \\ \hline
        Descripción general     & El sistema deberá permitir a los usuarios con rol de supervisor visualizar la información de producción correspondiente al día y turno seleccionados. \\ \hline
        Prioridad               & Media                                                                                                                                                 \\ \hline
    \end{tabularx}
\end{table}


\textbf{Programación de órdenes de producción}: El sistema deberá permitir a los usuarios con rol de programador cargar la programación de las órdenes de producción destinadas a los distintos centros productivos de la empresa (Tabla \ref{tab:req_ordenes}).

\begin{table}[H]
    \centering
    \caption{Requerimiento funcional - Programación de órdenes de producción}
    \label{tab:req_ordenes}
    \begin{tabularx}{\textwidth}{|X|X|}
        \hline
        Código de requerimiento & RF-07                                                                                                                                                                                                       \\ \hline
        Nombre                  & Programación de órdenes de producción                                                                                                                                                                       \\ \hline
        Propósito               & Gestionar la programación de órdenes destinadas a los centros de producción.                                                                                                                                \\ \hline
        Descripción general     & El sistema deberá permitir a los usuarios con rol de programador cargar y administrar la programación de las órdenes de producción que serán ejecutadas en los distintos centros productivos de la empresa. \\ \hline
        Prioridad               & Alta                                                                                                                                                                                                        \\ \hline
    \end{tabularx}
\end{table}


\textbf{Modificación de productos}: El sistema deberá permitir a los usuarios con rol de administrador modificar las especificaciones de los productos registradas en el Sistema Integral de Producción (Tabla \ref{tab:req_productos}).

\begin{table}[H]
    \centering
    \caption{Requerimiento funcional - Modificación de productos}
    \label{tab:req_productos}
    \begin{tabularx}{\textwidth}{|X|X|}
        \hline
        Código de requerimiento & RF-08                                                                                                                                                                \\ \hline
        Nombre                  & Modificación de productos                                                                                                                                            \\ \hline
        Propósito               & Mantener actualizadas las especificaciones de los productos registrados en el sistema.                                                                               \\ \hline
        Descripción general     & El sistema deberá permitir a los usuarios con rol de administrador modificar las especificaciones de los productos registradas en el Sistema Integral de Producción. \\ \hline
        Prioridad               & Media                                                                                                                                                                \\ \hline
    \end{tabularx}
\end{table}